\chapter{The First Day of the Royalty of Louis  XIV}
\lettrine{I}{n} the morning, the news of the death of the cardinal  was spread through the castle, and thence speedily reached the city. The  ministers Fouquet, Lyonne, and Letellier entered la salle des seances, to hold  a council. The king sent for them immediately. <Messieurs,> said  he, <as long as monsieur le cardinal lived, I allowed him to govern my  affairs; but now I mean to govern them myself. You will give me your advice  when I ask it. You may go.>

The ministers looked at each other with surprise. If they concealed a smile it  was with a great effort, for they knew that the prince, brought up in absolute  ignorance of business, by this took upon himself a burden much too heavy for  his strength. Fouquet took leave of his colleagues upon the stairs,  saying:—<Messieurs! there will be so much the less labour for  us.>

And he gayly climbed into his carriage. The others, a little uneasy at the turn  things had taken, went back to Paris together. Towards ten o'clock the  king repaired to the apartment of his mother, with whom he had a long and  private conversation. After dinner, he got into his carriage, and went straight  to the Louvre. There he received much company, and took a degree of pleasure in  remarking the hesitation of each, and the curiosity of all. Towards evening he  ordered the doors of the Louvre to be closed, with the exception of only one,  which opened on the quay. He placed on duty at this point two hundred Swiss,  who did not speak a word of French, with orders to admit all who carried  packages, but no others; and by no means to allow any one to go out. At eleven  o'clock precisely, he heard the rolling of a heavy carriage under the  arch, then of another, then of a third; after which the gate grated upon its  hinges to be closed. Soon after, somebody scratched with his nail at the door  of the cabinet. The king opened it himself, and beheld Colbert, whose first  word was this:—<The money is in your majesty's cellar.>

The king then descended and went himself to see the barrels of specie, in gold  and silver, which, under the direction of Colbert, four men had just rolled  into a cellar of which the king had given Colbert the key in the morning. This  review completed, Louis returned to his apartments, followed by Colbert, who  had not apparently warmed with one ray of personal satisfaction.

<Monsieur,> said the king, <what do you wish that I should  give you, as a recompense for this devotedness and probity?>

<Absolutely nothing, sire.>

<How! nothing? Not even an opportunity of serving me?>

<If your majesty were not to furnish me with that opportunity, I should  not the less serve you. It is impossible for me not to be the best servant of  the king.>

<You shall be intendant of the finances, M. Colbert.>

<But there is already a superintendent, sire.>

<I know that.>

<Sire, the superintendent of the finances is the most powerful man in the  kingdom.>

<Ah!> cried Louis, colouring, <do you think so?>

<He will crush me in a week, sire. Your majesty gives me a controle for  which strength is indispensable. An intendant under a  superintendent,—that is inferiority.>

<You want support—you do not reckon upon me?>

<I had the honour of telling your majesty, that during the lifetime of M.  de Mazarin, M. Fouquet was the second man in the kingdom; now M. de Mazarin is  dead, M. Fouquet is become the first.>

<Monsieur, I agree to what you told me of all things up to to-day; but  to-morrow, please to remember, I shall no longer suffer it.>

<Then I shall be of no use to your majesty?>

<You are already, since you fear to compromise yourself in serving  me.>

<I only fear to be placed so that I cannot serve your majesty.>

<What do you wish, then?>

<I wish your majesty to allow me assistance in the labours of the office  of intendant.>

<That post would lose its value.>

<It would gain in security.>

<Choose your colleagues.>

<Messieurs Breteuil, Marin, Hervart.>

<To-morrow the ordonnance shall appear.>

<Sire, I thank you.>

<Is that all you ask?>

<No, sire, one thing more.>

<What is that?>

<Allow me to compose a chamber of justice.>

<What would this chamber of justice do?>

<Try the farmers-general and contractors, who, during ten years, have  been robbing the state.>

<Well, but what would you do with them?>

<Hang two or three, and that would make the rest disgorge.>

<I cannot commence my reign with executions, Monsieur Colbert.>

<On the contrary, sire, you had better, in order not to have to end with  them.>

The king made no reply. <Does your majesty consent?> said Colbert.

<I will reflect upon it, monsieur.>

<It will be too late when reflection may be made.>

<Why?>

<Because you have to deal with people stronger than ourselves, if they  are warned.>

<Compose that chamber of justice, monsieur.>

<I will, sire.>

<Is that all?>

<No, sire; there is still another important affair. What rights does your  majesty attach to this office of intendant?>

<Well—I do not know—the customary ones.>

<Sire, I desire that this office be invested with the right of reading  the correspondence with England.>

<Impossible, monsieur, for that correspondence is kept from the council;  monsieur le cardinal himself carried it on.>

<I thought your majesty had this morning declared that there should no  longer be a council?>

<Yes, I said so.>

<Let your majesty then have the goodness to read all the letters  yourself, particularly those from England; I hold strongly to this  article.>

<Monsieur, you shall have that correspondence, and render me an account  of it.>

<Now, sire, what shall I do with respect to the finances?>

<Everything M. Fouquet has not done.>

<That is all I ask of your majesty. Thanks, sire, I depart in  peace;> and at these words he took his leave. Louis watched his  departure. Colbert was not yet a hundred paces from the Louvre when the king  received a courier from England. After having looked at and examined the  envelope, the king broke the seal precipitately, and found a letter from  Charles II. The following is what the English prince wrote to his royal  brother:—

<Your majesty must be rendered very uneasy by the illness of M. le  Cardinal Mazarin; but the excess of danger can only prove of service to you.  The cardinal is given over by his physician. I thank you for the gracious reply  you have made to my communication touching the Princess Henrietta, my sister,  and, in a week, the princess and her court will set out for Paris. It is  gratifying to me to acknowledge the fraternal friendship you have evinced  towards me, and to call you, more justly than ever, my brother. It is  gratifying to me, above everything, to prove to your majesty how much I am  interested in all that may please you. You are wrong in having Belle-Île-en-Mer  secretly fortified. That is wrong. We shall never be at war against each other.  That measure does not make me uneasy, it makes me sad. You are spending useless  millions; tell your ministers so; and rest assured that I am well informed;  render me the same service, my brother, if occasion offers.>

The king rang his bell violently, and his valet de chambre appeared.  <Monsieur Colbert is just gone; he cannot be far off. Let him be called  back!> exclaimed he.

The valet was about to execute the order, when the king stopped him.

<No,> said he, <no; I see the whole scheme of that man.  Belle-Isle belongs to M. Fouquet; Belle-Isle is being fortified: that is a  conspiracy on the part of M. Fouquet. The discovery of that conspiracy is the  ruin of the superintendent, and that discovery is the result of the  correspondence with England: this is why Colbert wished to have that  correspondence. Oh! but I cannot place all my dependence upon that man; he has  a good head, but I must have an arm!> Louis, all at once, uttered a  joyful cry. <I had,> said he, <a lieutenant of  musketeers!>

<Yes, sire—Monsieur d'Artagnan.>

<He quitted the service for a time.>

<Yes, sire.>

<Let him be found, and be here to-morrow the first thing in the  morning.>

The valet de chambre bowed and went out.

<Thirteen millions in my cellar,> said the king; <Colbert  carrying my purse and D'Artagnan my sword—I am king.>  